\documentclass[10pt,a4paper]{amsart}
\usepackage[margin=19mm]{geometry}
\usepackage{amsmath,amssymb,mathtools}
\usepackage[scr=rsfs,cal=euler]{mathalfa}
\usepackage{xfrac}
\usepackage[unicode,hidelinks,bookmarks=false]{hyperref}
\newcommand{\ie}{i.e.\ }
\newcommand{\cf}{cf.\ }
\newcommand{\resp}{respectively\ }
\newcommand{\Subsection}{\subsection}
\providecommand{\nobreakdash}{\nobreak\hbox{-}}
\setlength{\emergencystretch}{2em}
\multlinegap0pt
\usepackage{lmodern,array,longtable}
\hypersetup{
 pdftitle={Coven--Meyerowitz T2 necessity through coprime stripe collapse},
 pdfauthor={Jitendra Prajapati},
 pdfsubject={Integer tilings and the Coven--Meyerowitz conditions},
 pdfkeywords={integer tilings, cyclotomic polynomials, Coven--Meyerowitz, Lean}
}
\newtheorem{theorem}{Theorem}
\newtheorem{lemma}[theorem]{Lemma}
\newtheorem{corollary}[theorem]{Corollary}
\newcommand{\Z}{\mathbb Z}
\newcommand{\Q}{\mathbb Q}
\newcommand{\F}{\mathbb F}
\newcommand{\PhiC}{\Phi}
\newcommand{\code}[1]{\texttt{\small #1}}
\begin{document}
\title[Coven-Meyerowitz T2 necessity through coprime stripe collaps]{Coven-Meyerowitz T2 necessity through coprime stripe collapse}
\author{Jitendra Prajapati}
\date{}
\maketitle
\noindent\emph{Source and licence.} Coven-Meyerowitz T2 necessity through coprime stripe collapse. Version arXiv:2609.22280v1 (CC BY 4.0). This material is distributed under the Creative Commons Attribution 4.0 International licence, http://creativecommons.org/licenses/by/4.0/, as recorded in the arXiv OAI licence metadata for this exact version and shown on the abstract page https://arxiv.org/abs/2609.22280v1. Full rights notice: licenses/arxiv-2609.22280-CC-BY-4.0.txt.\par\smallskip
\noindent\textit{Editorial scope.} Counted unit: the original \texttt{theorem} environment \texttt{thm:cyclic} at source lines 199--201 together with its complete proof at lines 202--231; the delivered document keeps source lines 84--232 so that every referenced label and every citation resolves inside it. Native editable source is the paper's own concatenated TeX; the AMS wrapper is editorial formatting only. No publisher class, upstream script or unrelated artwork is redistributed.
\begin{equation}\label{eq:tiling}
 AB=J_N.
\end{equation}
The symbol $J_N$ is not the multiplicative identity. Use canonical residues to form ordinary mask polynomials. Evaluation at all $N$th roots separates group-ring elements by finite Fourier inversion.

\begin{lemma}[Exact primary allocation]\label{lem:allocation}
For an actual cyclic tiling $A+B=\Z_N$, the sets $S_A$ and $S_B$ partition the prime powers dividing $N$. If $q$ is prime, the number of $q$-primary levels owned by each factor is respectively $v_q(|A|)$ and $v_q(|B|)$. Neither factor has a primary cyclotomic divisor whose order does not divide $N$.
\end{lemma}
\begin{proof}
For $1\le e\le v_q(N)$, evaluation of~\eqref{eq:tiling} at a primitive $q^e$ root makes at least one factor vanish. An integer polynomial vanishing at such a root is divisible by its monic cyclotomic polynomial over $\Z$. Distinct cyclotomic divisors multiply to a divisor. Evaluating that product at $1$, and using $\PhiC_{q^e}(1)=q$, shows that the numbers $n_A,n_B$ of owned levels satisfy
\[
 v_q(N)\le n_A+n_B\le v_q(|A|)+v_q(|B|)=v_q(N).
\]
All inequalities are equalities, so each level is owned exactly once and both cardinality valuations are exhausted. Any further $q$-primary divisor would exceed this valuation. For $q\nmid N$, both cardinalities have $q$-valuation zero and cannot admit such a divisor.
\end{proof}

In particular, every T2 candidate order of either factor divides the period. The allocation argument is standard CM background~\cite[Lemma 2.1]{CM}; it does not by itself establish mixed zeros.

\section{Prime fibers, dephasing, and character identities}
Fix a prime $p\mid N$, write $N=pM$, and interchange the factors if necessary so that $B$ owns $\PhiC_p$. For $0\le i,j<p$ put
\[
 A_i=\{a\in\Z_M:i+pa\in A\},\qquad
 B_j=\{b\in\Z_M:j+pb\in B\}.
\]
The $B_j$ have the same positive cardinality $|B|/p$. Indeed, their count polynomial has degree at most $p-1$ and vanishes at a primitive $p$th root, so it is a scalar multiple of $1+z+\cdots+z^{p-1}$. The $A_i$ cannot all have the same cardinality, by exclusive ownership of $\PhiC_p$.

Write $M=KR$, where $K$ is the full $p$-primary part of $M$ and $(p,R)=1$. Choose $\rho$ with $p\rho\equiv1\pmod R$; if $R=1$, take $\rho=0$. In $\Z[\Z_M]$ define
\[
 U_i=T^{\rho i}A_i,\qquad V_j=T^{\rho j}B_j,\qquad F_{ij}=U_iV_j.
\]
The original cyclic carry is exactly retained by the identity
\begin{equation}\label{eq:carry}
 \left(\sum_i z^iA_i(T)\right)\left(\sum_j z^jB_j(T)\right)
 =J_M(T)\sum_{r=0}^{p-1}z^r
 \quad\text{modulo }(T^M-1,z^p-T).
\end{equation}
The quotient is identified with $\Z[X]/(X^{pM}-1)$ by $z=X$, $T=X^p$. It is not a replacement of the cyclic group by $\Z_p\times\Z_M$.

\begin{lemma}\label{lem:products}
The masks $F_{ij}$ are Boolean, invariant under translation by $R$, and satisfy
\begin{equation}\label{eq:rectangles}
 F_{ij}+F_{kl}=F_{il}+F_{kj},\qquad
 \sum_{i,j}F_{ij}=pJ_M.
\end{equation}
\end{lemma}
\begin{proof}
Let $\theta$ be an $M$th root of order $d>1$. If $p\mid d$, a $p$th root $\zeta$ of $\theta$ has order $pd$, and
\[
 [\Q(\zeta):\Q(\theta)]=\varphi(pd)/\varphi(d)=p.
\]
Thus $z^p-\theta$ is irreducible over $\Q(\theta)$. At $T=\theta$, the right side of~\eqref{eq:carry} vanishes. Each of the two fiber polynomials has degree below $p$, so irreducibility forces one entire coefficient family to vanish. Consequently every $F_{ij}(\theta)=0$.

If $(p,d)=1$, then $d\mid R$. Set $\lambda=\theta^\rho$, so $\lambda^p=\theta$ and $\lambda$ has order $d$. Substituting $z=\lambda w$ in~\eqref{eq:carry} makes
\[
 \left(\sum_i U_i(\theta)w^i\right)
 \left(\sum_j V_j(\theta)w^j\right)
 \quad\text{divisible by }w^p-1.
\]
Over $\Q(\theta)$, $\PhiC_p(w)$ is irreducible of degree $p-1$. Unless one factor is zero, one is a scalar multiple of $\PhiC_p$ and hence has a constant coefficient family. It follows that every rectangular difference vanishes at $\theta$. Evaluation at $w=1$ gives the vanishing of the total product there. The corresponding original root $\lambda$ is nontrivial because $d>1$.

At the trivial character the $V_j$ masses are equal, so the rectangular differences again vanish. The total mass is $|A||B|=pM$. Fourier inversion proves~\eqref{eq:rectangles}. For a fixed pair $i,j$, two representations in $A_i+B_j$ would give two representations in the original tiling with the same digit carry; hence $A_iB_j$, and its translate $F_{ij}$, are Boolean. Their Fourier coefficients vanish at every order containing $p$; the surviving orders divide $R$, which is exactly invariance under translation by $R$.
\end{proof}

This includes $p=2$. If $R=1$ there are no nontrivial $p$-free characters; if $M=1$ only the trivial-character and fixed-pair arguments are needed. Empty $A_i$ are allowed. Crucially, the periodicity conclusion concerns the \emph{products} $U_iV_j$, not the individual factors.

\section{Boolean stripes and coprime collapse}
At a fixed $x\in\Z_M$, the matrix $(F_{ij}(x))$ is Boolean, has zero additive rectangular differences, and contains exactly $p$ ones. Such a matrix is one complete row or one complete column of ones. To see this, if a row has unequal entries in two columns, the same nonzero difference occurs in every row; Booleanity forces those columns, and then all columns, to be constant. If no row varies, every row is constant. The total $p$ leaves exactly one stripe.

Consequently there are pairwise disjoint stripe sets $X_i,Y_j$ partitioning $\Z_M$ such that
\begin{equation}\label{eq:stripes}
 F_{ij}=X_i+Y_j.
\end{equation}
Each stripe mask is $R$-periodic, since it is uniquely decoded from the periodic matrix. Commutativity also gives the convolution minors
\begin{equation}\label{eq:minors}
 F_{ij}F_{kl}=F_{il}F_{kj}.
\end{equation}

\begin{lemma}[Coprime stripe collapse]\label{lem:collapse}
Let $p$ be prime and let $G$ be a finite abelian group with $p\nmid|G|$. Suppose Boolean stripe masks $X_0,\ldots,X_{p-1},Y_0,\ldots,Y_{p-1}$ partition $G$ and $F_{ij}=X_i+Y_j$ satisfy~\eqref{eq:minors}. Then all $X_i$ are empty or all $Y_j$ are empty.
\end{lemma}
\begin{proof}
Expand the minors to obtain $(X_i-X_k)(Y_j-Y_l)=0$. With $X=\sum_iX_i$ and $Y=\sum_jY_j$, sum over $k,l$ to get
\[
 (pX_i-X)(pY_j-Y)=0.
\]
Reducing in the commutative group algebra $\F_p[G]$ gives $XY=0$. Since $X+Y=J_G$ and $XJ_G=|X|J_G$, we get $X^2=|X|J_G$, and therefore
\[
 X^p=|X|^{p-1}J_G.
\]
Frobenius, however, gives $X^p=\sum_{x\in X}[px]$. Multiplication by $p$ permutes $G$, so the latter is a permutation of the original Boolean coefficient mask. A Boolean mask constant modulo $p$ is constant as an integer mask. Hence $X$ is empty or all of $G$, as required.
\end{proof}

To apply the lemma here, first descend the $R$-periodic masks from $\Z_M$ to $\Z_R$. For the lift $L$ of masks along reduction modulo $R$, a direct count of the $K=M/R$ lifts of a residue gives
\begin{equation}\label{eq:lift}
 L(f)L(g)=K L(fg).
\end{equation}
Equation~\eqref{eq:minors} therefore descends after cancellation of the nonzero integer $K$. This cancellation occurs in the integral group ring \emph{before} reduction modulo $p$; no division by $K$ in characteristic $p$ is used. The quotient stripes remain a Boolean partition and $(p,R)=1$, so Lemma~\ref{lem:collapse} applies.

Column-only would make $U_iV_j=Y_j$ independent of $i$. Taking masses and using the common positive mass of $V_j$ would make every $|A_i|$ equal, contradicting primary allocation. Thus the orientation is row-only:
\begin{equation}\label{eq:rowonly}
 U_iV_j=X_i,\qquad \sum_iX_i=J_M.
\end{equation}

\section{Actual lower tilings for every phase}\label{sec:phases}
For an arbitrary integer phase vector $h=(h_0,\ldots,h_{p-1})$, put
\[
 C_h=\sum_i T^{Rh_i}U_i.
\]
By~\eqref{eq:rowonly} and product periodicity,
\begin{equation}\label{eq:phase}
 C_hV_j=J_M\qquad\text{for all }j,h.
\end{equation}
The mask $C_h$ has nonnegative integer coefficients. Since $V_j$ is a nonempty Boolean mask, a coefficient of $C_h$ at least two would produce a convolution coefficient at least two, contradicting~\eqref{eq:phase}. Thus $C_h$ is Boolean and~\eqref{eq:phase} is an actual tiling of the strictly smaller period $M$. Translating the complement back also yields $C_hB_j=J_M$. We will use all phases, not just a single chosen quotient.

\section{Strong induction for both original factors}

\begin{theorem}\label{thm:cyclic}
For every $N\ge1$, both factors in an actual cyclic tiling $A+B=\Z_N$ satisfy T2.
\end{theorem}

\begin{proof}
Use strong induction on $N$. At $N=1$ both masks are the singleton identity and have no primary zeros. For $N>1$, choose $p\mid N$ and the labeling above. The induction hypothesis applies to every tiling~\eqref{eq:phase}.

\paragraph{The original factor $A$.}
For a $p$-free root $\theta$ whose order divides $M$, dephasing gives
\begin{equation}\label{eq:pfree}
 C_h(\theta)=\sum_i\theta^{\rho i}A_i(\theta)=A(\theta^\rho),
\end{equation}
since $\theta^R=1$ and $(\theta^\rho)^p=\theta$. Every selected $p$-free primary zero of $A$ thus passes to every $C_h$. If the target order $d$ is $p$-free, lower T2 for $C_0$ and~\eqref{eq:pfree}, with $\theta=\zeta^p$ for a primitive original root $\zeta$, give $A(\zeta)=0$.

If the selected order is $p^kd$ with $(p,d)=1$, then $k\ge2$, because $A$ does not own $\PhiC_p$. The identity
\[
 \PhiC_{p^k}(X)=\PhiC_{p^{k-1}}(X^p)
\]
and monic division by exponent residues modulo $p$ show that every $A_i$ is divisible by $\PhiC_{p^{k-1}}$. Explicitly, write the integer quotient of $A(X)$ by the displayed monic divisor as $\sum_iX^iQ_i(X^p)$ and compare residue classes. Thus every $C_h$ inherits the lower $p$-primary zero as well as the selected $p$-free zeros.

For a primitive root $\zeta$ of order $p^kd$, let $\theta=\zeta^p$, of order $p^{k-1}d\mid M$. Lower T2 gives $C_h(\theta)=0$ for every $h$. Comparing the zero phase with the phase changing only coordinate $i$ by one yields
\[
 \theta^{\rho i}(\theta^R-1)A_i(\theta)=0.
\]
Here $\theta^R$ has order $p^{k-1}>1$, so every $A_i(\theta)=0$. Regrouping gives $A(\zeta)=\sum_i\zeta^i A_i(\zeta^p)=0$. This is recovery of the original upper-$p$ mixed zero, not a claim that a single quotient preserves it.

\paragraph{The original factor $B$.}
We use the common-complement inheritance and lifting argument of Coven and Meyerowitz~\cite[Lemma~2.5]{CM}.
At zero phase, $C_0B_j=J_M$ for every $j$. Hence the dilated set $pC_0=\{pc:c\in C_0\}\subseteq\Z_{pM}$ is an actual complement of original $B$. Here dilation has mask $C_0(X^p)$; it is not scalar multiplication of coefficients by $p$. All lower $B_j$ have T2 by induction and share the same primary allocation, since their actual complement is the same $C_0$.

Every $p$-free primary zero of $B$ belongs to every $B_j$: compare exact allocation in $B+pC_0=\Z_{pM}$ with that in $B_j+C_0=\Z_M$. At a $p$-free prime-power order, the polynomial of $pC_0$ vanishes precisely when that of $C_0$ does, because raising a primitive root to its $p$th power preserves its order. Allocation then forces the zero into $B_j$.

If $B$ has a higher primary zero $p^k$, $k\ge2$, residue grading likewise puts $p^{k-1}$ in every $B_j$. For a $p$-free target $d$, lower T2 makes every $B_j(\zeta^p)$ zero and hence $B(\zeta)=0$. For a target $p^kd$ with $k\ge2$, apply the same argument at the lower target $p^{k-1}d$. For $k=1,d>1$, use the inherited lower target $d$. The remaining case $k=1,d=1$ is the already known $\PhiC_p$ zero. These cases exhaust all T2 requirements of both original factors by Lemma~\ref{lem:allocation}.
\end{proof}

\begin{thebibliography}{99}
\bibitem[CM]{CM} E. M. Coven and A. Meyerowitz, \emph{Tiling the integers with translates of one finite set}, Journal of Algebra \textbf{212} (1999), no.~1, 161--174. \url{https://arxiv.org/abs/math/9802122}.
\end{thebibliography}
\end{document}
